\clearpage\begingroup\SourceLeaf{54b}\HistoricalHead{49}{R}{ENCKE AN GAUSS.}
\begin{center}
[3.]\\[.7em]
\textsc{Encke} an \textsc{Gauss}. Berlin, 24. Mai 1842.
\end{center}
\noindent .\ .\ .\ .\ .\ .\ .\ .\ .\ .\ .\ .\ .\ .\ .\ .\ .\ .\ .\ .\ .
\par\noindent\hspace*{\parindent}\{Vor einigen Monaten erhielt ich von dem Chef der Artillerie eine Aufforderung
mich über die Art zu äussern, wie die Versuchsresultate der Artillerie
zu behandeln seyen, um sowohl einestheils Schusstafeln zu erhalten, welche
auch die möglichen Abweichungen geben, anderentheils bei künftigen Versuchen
eine Leitung zu haben und die Theorie zu vervollkommnen. Es war
dabei ausdrücklich bemerkt, dass man wünsche, ich solle von der bisherigen
ballistischen Theorie wo möglich keinen Gebrauch machen, da sie sich ungenügend
erwiesen. Als ein Beispiel waren folgende Zahlen gegeben, welche,
wie ich später erfuhr, zu einem 50 pf. Mörser, wo die Bombe etwa \(8\frac12\) Zoll
im Durchmesser hatte und das Geschoss etwa 6 Caliber lang war, gehörte.
\begin{center}\small\setlength{\tabcolsep}{3pt}\renewcommand{\arraystretch}{1.15}
\begin{tabular}{c@{\quad}c@{}lrrrrr}
&&&\multicolumn{5}{c}{Elevation}\\
Ladung && \multicolumn{1}{c}{(in Schritten)} & \(1^0\) & \(5^0\) & \(10^0\) & \(15^0\) & \(20^0\)\\
&\smash{\raisebox{-\normalbaselineskip}{\(\left\{\rule{0pt}{3\normalbaselineskip}\right.\)}}& Mittl. Schussweite & \(91{,}2\)&\(290{,}9\)&\(510{,}2\)&\(707{,}6\)&\(864{,}0\)\\
2 {\fontspec{FreeSerif.otf}℔} && »\quad Längenabweichung & \(4{,}5\)&\(14{,}7\)&\(14{,}5\)&\(21{,}1\)&\(23{,}5\)\\
&& »\quad Seitenabweichung & \(0{,}5\)&\(1{,}0\)&\(1{,}9\)&\(3{,}7\)&\(3{,}9\)\\[.6em]
&\smash{\raisebox{-\normalbaselineskip}{\(\left\{\rule{0pt}{3\normalbaselineskip}\right.\)}}& Mittl. Schussweite & \(201{,}7\)&\(754{,}2\)&\(1303{,}1\)&\(1729{,}9\)&\(2212{,}0\)\\
5 {\fontspec{FreeSerif.otf}℔} && »\quad Längenabweichung & \(14{,}8\)&\(22{,}7\)&\(34{,}3\)&\(26{,}5\)&\(41{,}4\)\\
&& »\quad Seitenabweichung & \(0{,}8\)&\(2{,}9\)&\(12{,}1\)&\(12{,}2\)&\(17{,}2\)\\[.6em]
&\smash{\raisebox{-\normalbaselineskip}{\(\left\{\rule{0pt}{3\normalbaselineskip}\right.\)}}& Mittl. Schussweite & \(358{,}0\)&\(1026{,}8\)&\(1716{,}8\)&\(2319{,}7\)&\(2832{,}4\)\\
{}[\textsuperscript{*})] && »\quad Längenabweichung & \(22{,}1\)&\(19{,}7\)&\(30{,}9\)&\(45{,}9\)&\(38{,}1\)\\
&& »\quad Seitenabweichung & \(0{,}9\)&\(5{,}0\)&\(11{,}3\)&\(17{,}0\)&\(31{,}9\)
\end{tabular}
\end{center}
\begingroup\renewcommand{\thefootnote}{}\footnotetext{[\textsuperscript{*}) Die Angabe der Ladung fehlt in der Handschrift.]}\endgroup
\HistoricalFoot{XI 1.}{7}
\clearpage\SourceLeaf{55}\HistoricalHead{50}{L}{PHYSIK. BRIEFWECHSEL.}
\noindent Bei jeder Elevation und Ladung wurden 15 Schuss gethan. In allem 225
Schuss.

Die mittlere Längenabweichung ist so zu verstehen, dass die Differenz
eines einzelnen Schusses vom Mittel aus allen ohne Rücksicht auf das Zeichen
genommen ist und daraus das Mittel gebildet. Ebenso die mittlere Seitenabweichung,
wo die wirkliche Richtung angenommen ist als mit der mittleren
Schusslinie gleichartiger Schüsse zusammenfallend.

Diese beobachteten Fehler schienen mir verhältnissmässig gering und sie
sind es auch wahrscheinlich nur dadurch geworden, dass auf die Lage des
Schwerpunktes des Geschosses genau Rücksicht genommen ist. Man bestimmt
jetzt die Lage des Schwerpunktes, indem man die Kugeln auf Quecksilber
schwimmen lässt. Wahrscheinlich hat man bei diesen Versuchen den Schwerpunkt
so gelegt, dass die Linie Schwerpunkt der äussern Figur und der
Masse mit der Axe des Rohrs zusammenfällt. Bei anderen Lagen ist die
Differenz der Schussweiten ungeheuer. Mir sind Versuche mitgetheilt, wo
unter ähnlichen Verhältnissen die Schussweite bei Schwerpunkt unten 838,0
Schritt betrug, bei Schwerpunkt oben aber 1362,1 Schritt.

Der Versuch, die bisher bekannten ballistischen Formeln mit diesen
Zahlen in Uebereinstimmung zu bringen, ist mir nicht geglückt. Wenn man
die Constante des Widerstandes und die Anfangsgeschwindigkeit aus irgend
zwei Werthen bestimmt, so weichen die übrigen viel zu stark ab. Ich habe
auch eine andere Hypothese versucht nämlich die Annahme, welche mir in
der Natur der Sache zu liegen scheint, dass bei dem Herausfahren der
Kugel aus der Mündung der Widerstand geringer ist als in den entfernteren
Theilen der Bahn, in so fern man annehmen kann, dass die Luft in gleichem
Sinne wie die Kugel in Bewegung gesetzt wird und sonach nicht als ein
ruhendes Medium zu betrachten sey. Indessen hat auch dieses nicht zu einem
erwünschten Ziele geführt.

Sonach habe ich durch die gewöhnlichen Interpolationsformeln die Bahn
für die zwischenliegenden Elevationsgrade ergänzt und, um von den Längenabweichungen
einigermassen Rechenschaft zu geben, die Differentialquotienten
in Bezug auf geänderte Elevation und Ladung numerisch aus den Interpolationsformeln
bestimmt und die mittlere Aenderung dieser beiden Elemente bestimmt,
welche hinreichen würde, um die Längenabweichungen wieder zu geben. In der
\clearpage\SourceLeaf{56}\HistoricalHead{51}{R}{ENCKE AN GAUSS.}
\noindent That findet sich dabei, dass man sehr nahe die Längenabweichungen erhält,
wenn man eine mittlere Variation von etwa 3 Minuten bei der Elevation um
\(\frac1{50}\) der Ladung annimmt. Die Seitenabweichungen sind fast ganz der Schussweite
proportional. Ausserdem habe ich empfohlen, sich womöglich an die
ballistischen Formeln anzuschliessen, um aus ihnen eine Form zu erhalten,
welche die Interpolation sicherer mache als die allgemeine Form der \textsc{Taylor}schen
Reihe.

Dieselbe Aufforderung war auch an \textsc{Bessel} ergangen und sein Gutachten
ist mir jetzt mitgetheilt worden. \textsc{Bessel} verwirft ganz alle Interpolation, so
lange man nicht aus der theoretischen Behandlung der Aufgabe eine bestimmte
Form erhalten hat, und spricht sich so aus, dass die grössten Fehler durch
die Interpolation auf dem gewöhnlichen Wege stattfinden würden, so dass man
bis zur Erhaltung der theoretischen Form jeden Versuch nur als ein einzelnes
isolirtes Faktum betrachten müsse.

Obgleich ich nun glaube, dass dieser Ausspruch vielleicht etwas scharf
abgefasst ist, um die Nothwendigkeit die Theorie zu verstehen eindringlich
zu machen, da \textsc{Bessel} doch auch als die möglicherweise anzuwendende Form,
wenn man sich diesen Fehlern aussetzen wolle, die gewöhnliche Interpolation
anführt und entwickelt, so hat mich doch der Ausspruch befremdet, weil ich
in der That nicht glaube, dass die Gefahr sehr grosser Fehler vorhanden
ist, wenn man die Versuche gehörig ausdehnt, auf die Intervalle an den
Grenzen des Anfangs und Endes weniger Gewicht legt und nicht ausserhalb
der Versuchswerthe hinausgeht. Namentlich ist mir aufgefallen, dass er sagt,
die Form der \textsc{Taylor}schen Reihe sei strenge nur bei unendlich kleinen Intervallen
anzuwenden.

Da wahrscheinlich der Gegenstand noch mehrfach gegen mich erwähnt
werden wird und da es mir in der That scheint, als wenn die Praxis hier der
Theorie vorgeeilt ist und die bisherigen Formeln nicht vermögen, innerhalb
der vorkommenden Fehlergrenzen die beobachteten Werthe wieder zu geben,
so möchte ich fast Sie ersuchen, Hochgeehrtester Herr Hofrath, diesem Gegenstande
einige Augenblicke zuzuwenden und mir gefälligst gelegentlich Ihre
Ansicht darüber mitzutheilen. Es werden grosse Summen auf die Versuche
gewendet und es wäre nicht unmöglich, dass, wenn eine befriedigende Theorie
\HistoricalFoot{}{7\textsuperscript{*}}
\clearpage\SourceLeaf{57}\HistoricalHead{}{L}{PHYSIK. BRIEFWECHSEL.}
\noindent aufgestellt würde, die Versuche so geordnet werden könnten, dass das Gesetz
des Widerstandes dadurch ein neues Licht erhielte.\quad .\ .\ .\ .\ .\ .\ .\ .\}
\clearpage\endgroup
